% Condition for cluster-approximated attention: objective, notation, the
% condition (cosh and Bennett variants), and the hybrid per-step certificate.
% Requires: tcolorbox, amsmath.

\section{Approximation}

\paragraph{Goal}
Suppose we split the KVs into clusters with close keys. Under what condition can we
faithfully approximate full attention? The condition must be efficiently verifiable,
and it must bound the approximation error.

\begin{tcolorbox}[colback=white, colframe=black, arc=0mm, boxrule=0.5pt]
For a query $q$ and a KV set clustered as $\{C_1,\dots,C_m\}$: for all $\epsilon > 0$,
using an approximation method $\mathbf{A}$ [efficient], under a condition
$\mathbf{Cond}$ [depending on $q$ and $\epsilon$, efficiently verifiable], the vanilla
output $o$ and the approximated output $\hat o$ satisfy
$\|\hat o - o\|_2 < \epsilon$.
\end{tcolorbox}

We use the \emph{average approximation}: each cluster is represented by its arithmetic
mean key and value $(\bar k_C, \bar v_C)$, $\bar k_C = \frac{1}{|C|}\sum_{i\in C} k_i$,
$\bar v_C = \frac{1}{|C|}\sum_{i\in C} v_i$ (subscripts dropped when clear).

\subsection{Notation}

\[
s_i = \frac{q k_i}{\sqrt{d_k}}, \quad
s_C = \frac{q \bar k}{\sqrt{d_k}}, \quad
Z_C = \sum_{i\in C} e^{s_i}, \quad
\hat Z_C = |C|\, e^{s_C}, \quad
u_C = \frac{\sum_{i\in C} e^{s_i} v_i}{Z_C},
\]
\[
p_C = \frac{Z_C}{\sum_{C'} Z_{C'}}, \qquad
\hat p_C = \frac{\hat Z_C}{\sum_{C'} \hat Z_{C'}}, \qquad
o = \sum_C p_C\, u_C, \qquad
\hat o = \sum_C \hat p_C\, \bar v_C .
\]
Per-cluster statistics (prompt-side summaries, $O(d_k)$ query-time work each):
\[
B \ge \max_i \|v_i\|_2, \qquad
B_C \ge \max_{i\in C} \|v_i\|_2, \qquad
\delta_C \ge \max_{i\in C} \Big|\frac{q(k_i - \bar k)}{\sqrt{d_k}}\Big|,
\]
\[
\sigma_C^2 \ \ge\ \frac{1}{|C|}\sum_{i\in C}(s_i - s_C)^2
\;=\; \frac{q\, \Sigma_C\, q^\top}{d_k},
\qquad \Sigma_C = \frac{1}{|C|}\sum_{i\in C}(k_i-\bar k)(k_i-\bar k)^\top .
\]

\emph{Remark.} Because $\bar k$ is the arithmetic mean,
$\sum_{i\in C}(s_i - s_C) = 0$ exactly; the mass-term analysis below relies on this
(a medoid or normalized centroid breaks it). It also gives $\sigma_C \le \delta_C$.

\subsection{Error decomposition}

Define
\[
b_C = \log \frac{Z_C}{\hat Z_C}
    = \log\Big(\frac{1}{|C|}\sum_{i\in C} e^{s_i - s_C}\Big) \ \ge\ 0
\quad\text{(Jensen)},
\qquad
A = \sum_{C'} \hat p_{C'}\, e^{b_{C'}} \ \ge\ 1,
\]
so that $p_C = \hat p_C\, e^{b_C} / A$. From
$\hat o - o = \sum_C (\hat p_C - p_C) u_C + \sum_C \hat p_C (\bar v_C - u_C)$
and $\|u_C\|_2 \le B$:
\[
\|\hat o - o\|_2
\ \le\
\underbrace{B \sum_C |p_C - \hat p_C|}_{\text{mass term}}
\; + \;
\underbrace{\sum_C \hat p_C\, \|\bar v_C - u_C\|_2}_{\text{value term}} .
\]

\subsection{Mass term}

Treating $e^{b_C}$ as a random variable with law $\hat p$ (mean $A$), and using
$\mathbb E|X - \mathbb E X| = 2\,\mathbb E (X - \mathbb E X)_+$:
\begin{align*}
\sum_C |p_C - \hat p_C|
  &= \frac{1}{A} \sum_C \hat p_C \big| e^{b_C} - A \big|
   = \frac{2}{A} \sum_C \hat p_C \big( e^{b_C} - A \big)_+ \\
  &\le \frac{2}{A} \sum_C \hat p_C \big( e^{b_C} - 1 \big)
   = \frac{2(A-1)}{A},
\end{align*}
using $A \ge 1$, $e^{b_C} \ge 1$. Since $A \mapsto 2(A-1)/A$ is increasing, it remains
to bound $e^{b_C}$.

\begin{tcolorbox}[colback=white, colframe=black, arc=0mm, boxrule=0.5pt]
\textbf{Lemma 1.} Let $x_i = s_i - s_C$, so $\frac{1}{|C|}\sum_i x_i = 0$,
$\frac{1}{|C|}\sum_i x_i^2 \le \sigma_C^2$, $x_i \le \delta_C$. Then
\[
e^{b_C} \;\le\; G(\sigma_C, \delta_C),
\qquad
G(\sigma,\delta) := \frac{\sigma^2 e^{\delta} + \delta^2 e^{-\sigma^2/\delta}}
                         {\sigma^2 + \delta^2} .
\]
\end{tcolorbox}

\emph{Proof.} $g(t) = (e^t - 1 - t)/t^2$ is nondecreasing, so for any $a$ and
$x \le \delta$,
\[
e^{x} = e^{a}\big[1 + (x-a) + g(x-a)(x-a)^2\big]
\ \le\ e^{a}\big[1 + (x-a) + g(\delta-a)(x-a)^2\big].
\]
Averaging over $i \in C$ (zero mean, second moment $\le \sigma^2$):
\[
e^{b_C} \ \le\ e^{a}\big[1 - a + g(\delta-a)\,(\sigma^2 + a^2)\big].
\]
Take $a = -\sigma^2/\delta$ and simplify. $\hfill\square$

\emph{Remarks.} $G(\delta,\delta) = \cosh\delta$: with no variance information the
lemma reduces to the chord bound $e^{b_C} \le \cosh\delta_C$, and $G(0,\delta) = 1$.
Only $x_i \le \delta_C$ is used, and the bound is tight. The proof holds verbatim with
over-estimated statistics $\hat\sigma_C \ge \sigma_C$, $\hat\delta_C \ge \delta_C$
(take $a = -\hat\sigma_C^2/\hat\delta_C$); one may always use
$\min\{\cosh\delta_C,\, G\}$, so the refined bound is never worse.

Writing $F_C$ for either $\cosh\delta_C$ or $G(\sigma_C,\delta_C)$:
\[
A \ \le\ S_F := \sum_C \hat p_C\, F_C ,
\qquad\qquad
\sum_C |p_C - \hat p_C| \ \le\ \frac{2\,(S_F - 1)}{S_F} .
\]

\subsection{Value term}

\begin{tcolorbox}[colback=white, colframe=black, arc=0mm, boxrule=0.5pt]
\textbf{Lemma 2.} Let $w_i = e^{s_i} / \sum_{j\in C} e^{s_j}$ with
$\max_{i\in C}|s_i - s_C| \le \delta_C$. Then
\[
\sum_{i\in C} \Big| \frac{1}{|C|} - w_i \Big| \ \le\ 2 \tanh(\delta_C/2).
\]
\end{tcolorbox}

\emph{Proof.} Let $Y_i = e^{s_i - s_C} \in [a, b]$, $a = e^{-\delta}$,
$b = e^{\delta}$, $T = \frac{1}{|C|}\sum_i Y_i$; the left side is
$\frac{1}{|C|}\sum_i |Y_i - T| / T$. Since $(y-T)_+$ is convex, it lies below its
chord on $[a,b]$: $(y-T)_+ \le (b-T)\frac{y-a}{b-a}$; averaging, and using
$\frac{1}{|C|}\sum_i |Y_i - T| = \frac{2}{|C|}\sum_i (Y_i - T)_+$,
\[
\sum_{i} \Big| \frac{1}{|C|} - w_i \Big|
\ \le\ \frac{2\,(b-T)(T-a)}{(b-a)\,T}
\ \le\ \frac{2\,(b-1)(1-a)}{b-a}
\ =\ 2\tanh(\delta/2),
\]
the middle maximization attained at $T = \sqrt{ab} = 1$. $\hfill\square$

Since $\sum_i (\frac{1}{|C|} - w_i) = 0$, for any fixed vector $c$,
\[
\bar v_C - u_C = \sum_{i\in C} \Big(\frac{1}{|C|} - w_i\Big) (v_i - c),
\qquad\text{hence}\qquad
\|\bar v_C - u_C\|_2 \ \le\ 2 B_C \tanh(\delta_C/2)
\]
with $c = 0$ (the centered choice $c = \bar v_C$ is also valid and never larger).

\subsection{The condition}

\begin{tcolorbox}[colback=white, colframe=black, arc=0mm, boxrule=0.5pt]
\begin{gather*}
\sum_{C} \hat p_{C}\left(
\frac{2B\,(F_C - 1)}{S_F}
+ 2 B_C \tanh(\delta_C/2)\right) \ <\ \epsilon
\quad\Longrightarrow\quad
\|\hat o - o\|_2 \ <\ \epsilon,
\\[4pt]
\text{where } S_F = \sum_{C} \hat p_C\, F_C,
\qquad F_C \in \{\cosh\delta_C,\ G(\sigma_C,\delta_C)\}.
\end{gather*}
\end{tcolorbox}

$F_C = \cosh\delta_C$ is the original condition; $F_C = G(\sigma_C,\delta_C)$ is
strictly tighter.

\subsection{Hybrid computation and per-step certificate}

At decode time a selected set $S$ of clusters is computed exactly \emph{inside the
same softmax}: for $C \in S$ the mass is the true $Z_C$ and the value is $u_C$ (so
$b_C = 0$), while $C \notin S$ keep $(\hat Z_C, \bar v_C)$. Let
$\hat o_{\mathrm{hyb}}$ denote this output; all $\hat p_C$ below are the
all-approximate weights, available before selection.

\begin{tcolorbox}[colback=white, colframe=black, arc=0mm, boxrule=0.5pt]
For any selected set $S$:
\begin{align*}
\|\hat o_{\mathrm{hyb}} - o\|_2
&\ \le\ \frac{2B\,T}{1+T}
 \;+\; \sum_{C\notin S} 2\,\hat p_C B_C \tanh(\delta_C/2),
\qquad T = \sum_{C \notin S} \hat p_C\, (F_C - 1),
\\[4pt]
&\ \le\ \sum_{C \notin S} \hat p_C
   \Big( 2B\,(F_C - 1) + 2 B_C \tanh(\delta_C/2) \Big).
\end{align*}
\end{tcolorbox}

\emph{Proof.} Run the same argument with the mixed estimators. For $C \in S$,
$e^{b_C} = 1$ and the value error vanishes, so $A \le 1 + \tilde T$ with
$\tilde T = \sum_{C\notin S} \hat p^{\mathrm{mix}}_C (F_C - 1)$, where
$\hat p^{\mathrm{mix}}$ uses the mixed normalizers. Since $Z_C \ge \hat Z_C$ on $S$,
the mixed normalizer is larger, hence $\hat p^{\mathrm{mix}}_C \le \hat p_C$ for
$C \notin S$ and $\tilde T \le T$; conclude by monotonicity of $x \mapsto x/(1+x)$.
The value term runs only over $C \notin S$. $\hfill\square$

The first bound is a per-decode-step error certificate computable from quantities the
selection pass already has; the second is additive per cluster and needs no
normalization across clusters.

\paragraph{Efficient verification}
Per cluster, store: $\bar k$, $\bar v$ (vectors); for $\delta_C$, either the coordinate
box $(k_{\max}, k_{\min})$ or the radius $r_C = \max_{i}\|k_i - \bar k\|_2$, giving
$\delta_C \le \|q\|_2\, r_C / \sqrt{d_k}$; for $\sigma_C^2$, the diagonal variance
$D_C = \operatorname{diag}(\Sigma_C)$ and the off-diagonal remainder
$\lambda_C = \|\Sigma_C - \operatorname{diag}(\Sigma_C)\|_2$, giving
\[
\sigma_C^2 \ \le\
\frac{1}{d_k}\Big( \sum_j q_j^2\, D_{C,j} + \lambda_C\, \|q\|_2^2 \Big);
\]
and the scalar $B_C$. All query-time work is $O(d_k)$ per cluster.
